\section{N Meetings in One Room}
\label{sec:n-meetings}

\problemstatement{We are given $n$ meetings, where the $i$-th meeting has a
start time $\textit{start}[i]$ and a finish time $\textit{end}[i]$. Only one
meeting room is available, so a meeting occupies the room for its entire
duration $[\textit{start}[i], \textit{end}[i]]$, and the room becomes free
only strictly after $\textit{end}[i]$ -- so two meetings can both be held in
the room only if one of them starts strictly after the other has ended. We
must select the \textbf{maximum number of non-overlapping meetings} that can
be held in this single room.}

\begin{patternbox}
\emph{``One resource (a room), many requests each with a start and an end,
maximize the number of non-conflicting requests you can serve''} is the
exact signature of the classical \emph{Activity Selection Problem}, the
flagship example of interval scheduling. The keyword combination to listen
for is \emph{``maximum number of non-overlapping intervals''} together with
\emph{a single shared resource}. The unintuitive but crucial insight that
makes this greedy rather than a search problem is that you should sort by
\textbf{finish time}, not by start time and not by duration -- a fact that
is worth memorizing because it differs from a tempting but incorrect first
instinct.
\end{patternbox}

\subsection*{Understanding the problem carefully}

It is tempting to think that scheduling the \emph{shortest} meetings first,
or the \emph{earliest-starting} meetings first, would leave the most room for
others. Both of these intuitions are wrong, and small examples expose the
flaw immediately. Consider two meetings, $A = [1, 100]$ and $B = [2, 3]$.
Sorting by start time would consider $A$ first (it starts earliest); taking
$A$ blocks $B$ entirely, leaving only $1$ meeting scheduled, whereas taking
$B$ instead would have left room afterward for further meetings starting
after time $3$. Sorting by duration would favor $B$ here by luck, but that
is not a reliable rule either -- a short meeting that starts very early and
a slightly longer meeting that finishes only one unit later can still flip
which choice is better, depending on what comes after.

What actually matters is not how long a meeting takes, nor when it starts,
but \emph{when it frees up the room for everyone else}. A meeting that
finishes early frees the room early, leaving the largest possible remaining
time window for future meetings -- regardless of how late it started or how
long it lasted.

\begin{ideabox}[Greedy Choice]
Sort all meetings by finish time, in increasing order. Always select the
next meeting (in this sorted order) whose start time is strictly after the
finish time of the most recently selected meeting.
\end{ideabox}

\subsection*{Why sorting by finish time is correct}

This is the standard exchange-argument proof used for activity selection.
Let $m_1$ be the meeting with the earliest finish time among all $n$
meetings. We claim there is always an optimal solution that includes $m_1$.
Suppose some optimal solution $O$ does not include $m_1$, but includes some
other meeting $m'$ as its first (earliest-finishing) chosen meeting instead.
Since $m_1$ has the globally earliest finish time of \emph{all} meetings, in
particular $\textit{end}[m_1] \le \textit{end}[m']$. Replacing $m'$ with
$m_1$ inside $O$ cannot create any new conflict with the meetings that come
after $m'$ in $O$, because $m_1$ frees the room \emph{no later} than $m'$
did. So this swap produces another solution that is still feasible and has
exactly the same number of meetings as $O$ -- meaning some optimal solution
does include $m_1$. Once $m_1$ is fixed as selected, the remaining problem
is exactly the same problem, restricted to meetings that start after
$\textit{end}[m_1]$, so the same argument applies recursively (or, in
iterative form, repeatedly) to the next-earliest-finishing remaining
meeting, and so on. This is what justifies scanning the finish-time-sorted
list exactly once, greedily picking whichever next meeting is compatible
with the last one picked.

\subsection*{Algorithm}

\begin{enumerate}
  \item Sort all $n$ meetings by $\textit{end}[i]$ in increasing order.
  \item Initialize $\textit{count} \gets 1$ and
        $\textit{lastEnd} \gets \textit{end}$ of the first meeting in
        sorted order (the very first meeting is always selected, since
        nothing has been scheduled yet to conflict with it).
  \item For each remaining meeting $i$ in sorted order, starting from the
        second one:
  \begin{enumerate}
    \item If $\textit{start}[i] > \textit{lastEnd}$: select this meeting.
          $\textit{count} \gets \textit{count} + 1$;
          $\textit{lastEnd} \gets \textit{end}[i]$.
    \item Otherwise: skip this meeting (it conflicts with the last
          selected meeting).
  \end{enumerate}
  \item Return \textit{count}.
\end{enumerate}

\begin{algorithm}[H]
\caption{Maximum Meetings in One Room}
\begin{algorithmic}[1]
\Procedure{MaxMeetings}{$\textit{start}, \textit{end}$}
  \State sort meetings by $\textit{end}[i]$ in increasing order
  \State $\textit{count} \gets 1$
  \State $\textit{lastEnd} \gets \textit{end}[0]$ \Comment{always take the first meeting}
  \For{$i \gets 1$ to $n-1$}
    \If{$\textit{start}[i] > \textit{lastEnd}$}
      \State $\textit{count} \gets \textit{count} + 1$
      \State $\textit{lastEnd} \gets \textit{end}[i]$
    \EndIf
  \EndFor
  \State \Return \textit{count}
\EndProcedure
\end{algorithmic}
\end{algorithm}

\subsection*{Dry run}

\dryrun{Take meetings with $(\textit{start}, \textit{end})$ pairs:
$(1,2), (3,4), (0,6), (5,7), (8,9), (5,9)$.}

Sorting by finish time gives the order:
\[
  (1,2),\ (3,4),\ (0,6),\ (5,7),\ (8,9),\ (5,9).
\]

\begin{itemize}
  \item Select $(1,2)$ unconditionally. $\textit{count}=1$,
        $\textit{lastEnd}=2$.
  \item $(3,4)$: $\textit{start}=3 > \textit{lastEnd}=2$. Select.
        $\textit{count}=2$, $\textit{lastEnd}=4$.
  \item $(0,6)$: $\textit{start}=0 \not> \textit{lastEnd}=4$. Skip.
  \item $(5,7)$: $\textit{start}=5 > \textit{lastEnd}=4$. Select.
        $\textit{count}=3$, $\textit{lastEnd}=7$.
  \item $(8,9)$: $\textit{start}=8 > \textit{lastEnd}=7$. Select.
        $\textit{count}=4$, $\textit{lastEnd}=9$.
  \item $(5,9)$: $\textit{start}=5 \not> \textit{lastEnd}=9$. Skip.
\end{itemize}

The maximum number of meetings that can be held is $\boxed{4}$, namely
$(1,2), (3,4), (5,7), (8,9)$.

\subsection*{C++ implementation}

\begin{lstlisting}
#include <bits/stdc++.h>
using namespace std;

int maxMeetings(vector<int>& start, vector<int>& end) {
    int n = start.size();
    vector<pair<int,int>> meetings(n);
    for (int i = 0; i < n; i++) {
        meetings[i] = {end[i], start[i]};
    }

    sort(meetings.begin(), meetings.end()); // sorts by end time first

    int count = 1;
    int lastEnd = meetings[0].first;

    for (int i = 1; i < n; i++) {
        int s = meetings[i].second;
        int e = meetings[i].first;
        if (s > lastEnd) {
            count++;
            lastEnd = e;
        }
    }
    return count;
}

int main() {
    vector<int> start = {1, 3, 0, 5, 8, 5};
    vector<int> end   = {2, 4, 6, 7, 9, 9};

    cout << "Maximum meetings: " << maxMeetings(start, end) << endl;
    return 0;
}
\end{lstlisting}

\begin{complexitybox}
\textbf{Time Complexity.} Sorting the $n$ meetings by finish time costs
$O(n \log n)$. The subsequent single linear pass costs $O(n)$. So the
overall time complexity is
\[
  O(n \log n).
\]
\textbf{Space Complexity.} We use an auxiliary array of $n$ pairs to sort
by finish time, giving $O(n)$ space (this can be reduced to $O(1)$
auxiliary space if the input is given as an array of pairs already, sorted
in place).
\end{complexitybox}

\begin{takeawaybox}
For interval-scheduling problems with a single shared resource, always sort
by \textbf{finish time}, never by start time or duration. The meeting that
frees the resource soonest is always at least as good a first choice as any
other meeting, by a direct exchange argument -- and this single fact is the
entire content of the activity selection problem, a pattern that reappears
throughout the rest of this chapter under different names.
\end{takeawaybox}
